% ---
% title: A combinatorial proof
% tags: combinatorial-proof, combinations
% ---

\begin{problem}
	Prove Pascal's rule:
	\[ \binom nk+\binom n{k+1}=\binom{n+1}{k+1} \qquad (0\le k<n), \]
	(a) by algebra, (b) by a counting argument.
\end{problem}

\begin{slides}
	\small
	\exstatement
	\pause
	\sol \textbf{Algebra.} Common denominator $(k+1)!\,(n-k)!$:
	\[
		\frac{n!}{k!(n-k)!}+\frac{n!}{(k+1)!(n-k-1)!}
		= \frac{n!\,\big[(k+1)+(n-k)\big]}{(k+1)!\,(n-k)!}
		= \frac{(n+1)!}{(k+1)!\,(n-k)!} = \binom{n+1}{k+1},
	\]
	since $n!\,(n+1)=(n+1)!$ and $n-k=(n+1)-(k+1)$.\\[1mm]
	\textbf{Story.} A bowl holds $n$ apples and $1$ pear. Choose $k+1$ fruits:
	$\binom{n+1}{k+1}$ ways. Split the choices by the pear:
	\emph{pear not chosen} (all $k+1$ are apples): $\binom n{k+1}$;
	\emph{pear chosen} ($k$ more apples from $n$): $\binom nk$.
	The same set counted in two ways $\Rightarrow$ the identity. This is the
	rule that builds Pascal's triangle.
\end{slides}

\begin{solution}
	\idea (a) Bring the two fractions to a common denominator. (b) Find a
	counting question whose answer is $\binom{n+1}{k+1}$, and answer it a second
	way, by splitting into two cases.

	\step{(a) Algebra.} Write both terms with the common denominator
	$(k+1)!\,(n-k)!$. Since $(k+1)!=(k+1)\,k!$ and $(n-k)!=(n-k)\,(n-k-1)!$:
	\[
		\frac{n!}{k!\,(n-k)!}+\frac{n!}{(k+1)!\,(n-k-1)!}
		=\frac{n!\,(k+1)}{(k+1)!\,(n-k)!}+\frac{n!\,(n-k)}{(k+1)!\,(n-k)!}
		=\frac{n!\,(n+1)}{(k+1)!\,(n-k)!}.
	\]
	Now $n!\,(n+1)=(n+1)!$ and $n-k=(n+1)-(k+1)$, so this equals
	\[ \frac{(n+1)!}{(k+1)!\,\big((n+1)-(k+1)\big)!}=\binom{n+1}{k+1}. \]

	\step{(b) Counting (the ``apple and pear'' argument).} A bowl contains $n$
	apples and $1$ pear, all different. Count the ways to choose $k+1$ fruits.
	\begin{itemize}[nosep]
		\item Directly: $k+1$ out of $n+1$ fruits: $\binom{n+1}{k+1}$.
		\item By whether the pear is chosen. \emph{Pear not chosen}: all $k+1$
		      fruits are apples, $\binom n{k+1}$ ways. \emph{Pear chosen}: the
		      other $k$ fruits are apples, $\binom nk$ ways. These two cases are
		      disjoint and cover everything, so the total is
		      $\binom nk+\binom n{k+1}$.
	\end{itemize}
	The same set was counted in two ways, so the two expressions are equal.
	\begin{answer}
		$\binom nk+\binom n{k+1}=\binom{n+1}{k+1}$
	\end{answer}
	\begin{remark}
		Example: $\binom41+\binom42=4+6=10=\binom52$. This rule builds Pascal's
		triangle; in the binomial tree (Exercise~\exref{binomial-tree}) it says
		that the paths into a node come from the two nodes feeding it.
		Counting proofs avoid long fraction manipulations; Vandermonde's identity
		(Exercise~\exref{vandermonde}) is another example.
	\end{remark}
\end{solution}
